\section{Recursive propagation of the original signed metric bounds}
\label{GraphV2:sec:recursive-metric-transport}

This proof combines the earlier full-relative-spectrum estimate AW16,
the exact trace norm and actual-overlap estimates of SP, the supplied
centered Hilbert--Schmidt estimate, and the all-degree angle estimate
EP17--EP25 on their common original source.  Each of those complete
proofs is retained in this edition.  The construction below carries
their pointwise minimum through both the original volume and the exact
nonlinear endpoint coordinate, without changing either endpoint form.

\subsection{The original source and all four quotient volumes}

Fix the original packet \(h\), tensor order \(k\geq1\), monic cyclic
polynomial \(\chi=\chi_{h,k}\) of degree \(q\geq1\), and
\[
 H=\mathcal P_{2q},\qquad S=k/2+iy,\qquad
 v(y)=(1,S,\ldots,S^{2q}).
 \tag{RMT1}
\]
The coordinates are the original ascending monomial coordinates.
The original reference and arithmetic measures are
\[
 d\mu_0=r_{1/4}^{*k}(y)\,dy,\qquad
 d\mu_1=w_h^{*k}(y)\,dy,\qquad
 w_h(y)=\frac{|(2\xi/h)(1/2+iy)|^2}{2\pi}.
 \tag{RMT2}
\]
All their original mass and Taylor-unit factors remain.  In particular
the exact reference convolution identity is
\[
 r_{1/4}^{*k}
 =\frac{c_{1/4}^{\,k}}{c_{k/4}}r_{k/4},\qquad
 c_a=2^{1-2a}\Gamma(2a),\qquad c_{1/4}=\sqrt{2\pi}.
 \tag{RMT3}
\]
Thus these are the same measures in AW and SP.  Every argument below
also applies to the actual finite-circle comparison when one of the
measures is its specified positive atomic measure on this same line
and its retained degree-\(2q\) Gram is positive definite.  That
positivity is already proved by the original sampled-source bounds.
This statement does not assign moment structure to an arbitrary
positive coefficient matrix.

For \(0\leq x\leq1\) let
\[
 M_x=\int v(y)^*v(y)\,d((1-x)\mu_0+x\mu_1)(y),\qquad
 \dot M=M_1-M_0,\qquad C_x=M_x^{-1}\dot M.
 \tag{RMT4}
\]
The two positive endpoint forms make \(M_x\) positive definite
throughout the closed interval.  For every original cutoff \(N\),
let \(I_N:\mathcal P_N\to H\) be coefficient inclusion and
\(B_N:\mathcal P_{N-q}\to H\) multiplication by \(\chi\), followed
by inclusion.  Negative relation degree means the zero space.
Retain the orthogonal projections
\[
 P_N=I_N(I_N^*M_xI_N)^{-1}I_N^*M_x,\qquad
 Q_N=B_N(B_N^*M_xB_N)^{-1}B_N^*M_x,\quad Q_{q-1}=0.
 \tag{RMT5}
\]
AW's relation matrix has target \(\mathcal P_N\), so its exact
relation to this one is \(B_N=I_N B_N^{\rm AW}\).
Equation (RMT5), after that substitution, is precisely AW4 and SP2.
The quotient map remains \(J_N:\mathcal P_N\to E=\mathbb{C}[S]/(\chi)\),
and its canonical positive quotient Gram is \(G_N(x)\).
Put \(V_N(x)=\det G_N(x)\) and retain the four signs:
\[
 \mathcal B(x)=\log\frac{V_{q-1}(x)V_q(x)}
                              {V_{2q-1}(x)V_{2q}(x)},\qquad
 U_x=Q_{2q-1}+Q_{2q}-Q_q,\qquad
 W_x=P_{2q-1}-P_{q-1}+P_{2q}-P_q.
 \tag{RMT6}
\]
The relation-first determinant-line formula and differentiation of
each of these four determinants give, without commuting their
operators,
\[
 \mathcal B'(x)=\operatorname{Tr}((U_x-W_x)C_x).
 \tag{RMT7}
\]
This is the same actual integrand as AW6 and SP7.

\subsection{The simultaneous bounds on that integrand}

Let \(b_1\leq\cdots\leq b_{2q+1}\) be every eigenvalue, with
multiplicity, of \(D=M_0^{-1}M_1\) in the original \(M_0\) metric.
They are positive.  Retain
\[
 c_j(x)=\frac{b_j-1}{1-x+xb_j},\qquad
 d_x=c_{2q+1}(x)-c_1(x),\qquad
 \kappa=\frac{b_{2q+1}}{b_1}.
 \tag{RMT8}
\]
These are the ordered eigenvalues of \(C_x\).  Indeed
\(C_x=((1-x)I+xD)^{-1}(D-I)\), and differentiation with
respect to \(b_j\) gives the positive quantity
\((1-x+xb_j)^{-2}\).  In particular
\(\int_0^1d_x\,dx=\log\kappa\).

Set \(s_x=\dim(\operatorname{ran}U_x\cap\operatorname{ran}W_x)\).
The two ranks are \(q+1\) in dimension \(2q+1\), so \(s_x\geq1\).
All traces and trace norms below are in the same original \(M_x\)
metric.  Define the two calculated functions
\[
 \begin{split}
 E_q(x)&=c_{q+2}(x)-c_q(x)
       +2\sum_{j=1}^{q-1}(c_{2q+2-j}(x)-c_j(x)),\\
 O_q(x)&=\frac{d_x}{2}\min\left\{
 4q-2s_x,\
 \sqrt{(2q+1)\bigl(8q-4-2\operatorname{Tr}(U_xW_x)\bigr)}
 \right\}.
 \end{split}
 \tag{RMT9}
\]
The sum is empty when \(q=1\).
The complete nested-flag calculation gives both \(U_x,W_x\)
the spectrum \(0^{[q]},1^{[2]},2^{[q-1]}\).  Decomposing each
as the sum of its spectral projections of ranks \(q+1\) and
\(q-1\), and applying the ordered projection-trace inequalities,
gives \(|\mathcal B'|\leq E_q\), exactly as in AW15--AW16.
For clarity the subtraction of the two bounds cancels the
common \(c_{q+1}\) only after both sums have been written; it
leaves every term displayed in \(E_q\).

Both positive operators dominate the original orthogonal projection
onto their common range.  Subtracting that same projection from
both leaves positive operators with traces \(2q-s_x\), whence
\(\|U_x-W_x\|_1\leq4q-2s_x\).  Also their computed spectra give
\[
 \operatorname{Tr}\bigl((U_x-W_x)^2\bigr)
 =8q-4-2\operatorname{Tr}(U_xW_x).
 \tag{RMT10}
\]
Cauchy--Schwarz on all \(2q+1\) eigenvalues of \(U_x-W_x\)
gives the second trace-norm bound in (RMT9).
Since \(\operatorname{Tr}(U_x-W_x)=0\), subtract
\((c_1+c_{2q+1})I/2\) from \(C_x\) in (RMT7).
In an \(M_x\)-orthonormal eigenbasis of \(U_x-W_x\), the
absolute value of each diagonal entry of this centered \(C_x\)
is at most \(d_x/2\).  Multiplying by the corresponding
absolute eigenvalues and summing gives the retained exact
trace-norm bound
\[
 |\mathcal B'|\leq T_q(x):=\frac{d_x}{2}\|U_x-W_x\|_1
                   \leq O_q(x).
 \tag{RMT10a}
\]
The second inequality follows from the two trace-norm
estimates just proved.  Thus neither of those majorants
replaces the original trace norm.

The user's additional supplied projector calculation gives another
bound on this same integrand.  Define
\[
 \Sigma_x=\operatorname{Tr}(C_x^2)
              -\frac{(\operatorname{Tr}C_x)^2}{2q+1},
 \qquad
 S_q(x)=\sqrt{\Sigma_x
                 \operatorname{Tr}\bigl((U_x-W_x)^2\bigr)}.
 \tag{RMT10b}
\]
The nonnegative number \(\Sigma_x\) is the squared
Hilbert--Schmidt norm of
\(C_x-\operatorname{Tr}(C_x)I/(2q+1)\) in the original
\(M_x\) Hilbert space.  The trace of \(U_x-W_x\) is zero,
so Hilbert--Schmidt Cauchy--Schwarz gives
\(|\mathcal B'|\leq S_q\).  Equivalently one may apply the
isometry \(M_x^{1/2}:(H,M_x)\to(H,I)\) simultaneously to
both operators; conjugation preserves their traces and
products and proves the same inequality.  This is a
calculation with the original metric, not a change of
either endpoint source form.

The original crossed quotient pairs are \((q-1,2q)\) and
\((q,2q-1)\).  Retain their full canonical minimum-section angle
lists \((\theta_1,\ldots,\theta_q)\) and
\((\varphi_1,\ldots,\varphi_q)\).  In the second list, ordered
with a smallest angle first, the original ambient dimension
forces \(\varphi_1=0\).  The count is proved in EP17--EP25.
For the following scalar sums put \(a=2q-1\) and define the
explicitly indexed list
\[
 \lambda_i=-\log\cos^2\theta_i\quad(1\leq i\leq q),\qquad
 \lambda_{q+j-1}=-\log\cos^2\varphi_j\quad(2\leq j\leq q).
\]
Then \(\mathcal B=-\log\cos^2\varphi_1+
\sum_{\nu=1}^{a}\lambda_\nu=0+\sum_{\nu=1}^{a}\lambda_\nu\).
Every additional zero remains in the \(\lambda\)-list, while
the designated geometric entry \(\varphi_1\) is still displayed
in the full second list.  The exact projection-difference
spectra give
\[
 |\mathcal B'(x)|
 \leq d_x\left(\sin\varphi_1+
               \sum_{\nu=1}^{a}\sqrt{1-e^{-\lambda_\nu(x)}}\right)
 \leq a\sqrt{1-e^{-\mathcal B(x)/a}}\,d_x.
 \tag{RMT11}
\]
The second inequality is concavity of \(z\mapsto\sqrt{1-e^{-z}}\)
on \([0,\infty)\), applied to the full \(a\)-term list, including
its zeros, together with the displayed \(\sin\varphi_1=0\).
Both \(q\)-entry lists remain the geometric data; the known
zero contribution is kept explicitly in these scalar identities.

The needed strict positivity of \(\mathcal B\) follows on the
actual measures rather than being an additional assumption.
If \(\mathcal B=0\), both nested quotient pairs have equal
minimum spaces.  To see every implication, their canonical
minimum sections satisfy
\(G_i-G_j=(R_i-R_j)^*M_x(R_i-R_j)\succeq0\).
Thus each determinant ratio in the crossed pairs is at least
one.  A zero sum of their nonnegative logarithms makes both
ratios one.  All eigenvalues of \(G_j^{-1/2}G_iG_j^{-1/2}\)
are at least one, so determinant one makes every eigenvalue
one and \(G_i=G_j\).  The displayed squared norm then gives
\(R_i=R_j\) as maps to the same original source \(H\).
In the first pair \(R_{q-1}\) is the literal degree-less-than-\(q\)
remainder section.  Its equality with \(R_{2q}\) puts the
constant polynomial \(1\) in the latter minimum space,
which is orthogonal to \(\ker J_{2q}=\chi\mathcal P_q\).
For \(\chi(S)=\sum a_jS^j\), put
\(\chi^{\#_k}(S)=\sum\overline{a_j}(k-S)^j\).
This belongs to \(\mathcal P_q\) and equals
\(\overline{\chi(S)}\) on \(S=k/2+iy\).  The asserted
orthogonality would therefore give
\[
 0=\int\chi(S)\chi^{\#_k}(S)\,d\mu_x
   =\int|\chi(S)|^2\,d\mu_x,
 \tag{RMT12}
\]
contradicting the positive Gram of the nonzero polynomial
\(\chi\).  This proof also covers the specified positive
atomic comparison.  Thus \(\mathcal B(x)>0\) on the entire
closed interval.

\subsection{Carrying the refined bound through both endpoint coordinates}

Define the positive function
\[
 z_x=a\sqrt{1-e^{-\mathcal B(x)/a}}>0.
\]
With every preceding control retained, set
\[
 \begin{split}
 A_q(x)&=z_xd_x,\\
 J_q&=\int_0^1\min\{E_q(x),T_q(x),O_q(x),S_q(x),A_q(x)\}\,dx,\\
 I_q&=\int_0^1\min\left\{d_x,\
 \frac{E_q(x)}{z_x},\frac{T_q(x)}{z_x},\
 \frac{O_q(x)}{z_x},\frac{S_q(x)}{z_x}\right\}\,dx .
 \end{split}
 \tag{RMT13}
\]
These are actual quantities of the original finite source path.
The projections are continuous matrix functions of \(x\).
Explicitly \(s_x=2(q+1)-\operatorname{rank}[\,U_x\ W_x\,]\).
A rank superlevel set is a finite union of sets where a
specified minor is nonzero, hence is Borel.  This proves
measurability of \(s_x\) and of every displayed integrand.
Equation (RMT12)
and compactness give a positive lower bound for \(\mathcal B\)
on this fixed path.  Hence all denominators in \(I_q\) are
nonzero, and the integrands are finite and bounded.  This is
a fixed-source statement and asserts no bound uniform in
the tensor order.

All five preceding estimates bound the same \(|\mathcal B'|\).
Pointwise comparison followed by the triangle inequality
for the integral therefore proves
\[
 \begin{gathered}
 |\mathcal B(1)-\mathcal B(0)|\leq J_q,\\
 J_q\leq
 \min\left\{\mathcal L_q(D),\int_0^1T_q(x)\,dx,
             \int_0^1O_q(x)\,dx,\int_0^1S_q(x)\,dx,
                         \int_0^1A_q(x)\,dx\right\}.
 \end{gathered}
 \tag{RMT14}
\]
where the earlier AW13 quantity is retained in full:
\[
 \mathcal L_q(D)=\log\frac{b_{q+2}}{b_q}
       +2\sum_{j=1}^{q-1}\log\frac{b_{2q+2-j}}{b_j}
       \leq(2q-1)\log\kappa.
 \tag{RMT15}
\]
Indeed integration of each \(c_j\) gives \(\log b_j\),
which proves the first entry of the minimum in (RMT14).

For \(B>0\) let
\[
 \begin{gathered}
 \mathscr H_a(B)=2\operatorname{arcosh}(e^{B/(2a)}),\qquad
 \mathscr H_a'(B)=
 \frac{1}{a\sqrt{1-e^{-B/a}}},\\
 \mathscr H_a^{-1}(z)=2a\log\cosh(z/2)\quad(z>0).
 \end{gathered}
 \tag{RMT16}
\]
The derivative follows by differentiating the positive real
inverse branch of \(\cosh\); substitution proves both inverse
composites.  Applying this exact derivative to each of the
five bounds on \(\mathcal B'\) gives
\[
 |\mathscr H_a(\mathcal B(1))-\mathscr H_a(\mathcal B(0))|
 \leq I_q\leq\int_0^1d_x\,dx=\log\kappa.
 \tag{RMT17}
\]
Thus the exact trace norm, supplied variance estimate, overlap
and full-relative-spectrum information propagate into the nonlinear bound,
not only into its linear consequence.

Writing \(H_0=\mathscr H_a(\mathcal B(0))\), the inverse
in (RMT16), extended continuously by value zero at zero,
and the additive estimate in (RMT14) give the simultaneous
interval
\[
 \begin{split}
 \mathcal B(1)&\geq
 \max\left\{0,\ \mathcal B(0)-J_q,\
   2a\log\cosh\left(\frac{\max\{0,H_0-I_q\}}2\right)\right\},\\
 \mathcal B(1)&\leq
 \min\left\{\mathcal B(0)+J_q,\
   2a\log\cosh\left(\frac{H_0+I_q}2\right)\right\}.
 \end{split}
 \tag{RMT18}
\]
The actual \(\mathcal B(1)\) is strictly positive although
the lower bound can equal zero.  If \(M_1=cM_0\) with
\(c>0\), all the \(c_j\) coincide, so \(d_x,E_q,T_q,O_q,S_q,A_q\),
\(I_q,J_q\) vanish.  Each original quotient determinant is
multiplied by the same power \(q\) of \(1-x+xc\);
the four signs then give \(\mathcal B(1)=\mathcal B(0)\),
as (RMT18) also states.  This includes \(q=1\), for which
the ordered sums above are empty and \(a=1\).

Every projection and metric in this argument uses the
original \(J_N\), \(B_N\), coefficient basis and source
measure.  The arithmetic class map
\(\sigma_h^{\otimes k}\eta_{h,k}J_N\), its complete Taylor
unit, and the original theta primitives of differences of
minimum sections are unchanged.  The proof refines numerical
bounds on those exact maps; it introduces no identification
between their source kernels or between distinct support
labels.  The earlier complete AW, SP and EP proofs, together
with the explicit correspondences (RMT3)--(RMT7), provide
the full input proofs used in this propagation.

\subsection{Carrying the constructed signed residual into the endpoint}

The complete supplied projector correspondence and its explicit
finite approximation are proved in
Section~\ref{GraphV2:sec:incoming-source-metric-control}, (ISM14)--(ISM37).
Their notation has the exact identities
\[
 X_x=C_x,\qquad \mathcal A_x=U_x-W_x,\qquad
 \Pi_N=P_N-Q_N.
\]
For completeness retain the actual finite construction here.
For each integer \(L\geq0\), put
\[
 \begin{aligned}
 B_x^\circ&=(1-x)I+xD,&
 \alpha_x&=\frac{2(1-x)+x(b_1+b_{2q+1})}{2},\\
 H_x^\circ&=I-B_x^\circ/\alpha_x,&
 C_x^{[L]}&=\alpha_x^{-1}\sum_{r=0}^{L}
                      (H_x^\circ)^r(D-I),\\
 R_x^{[L]}&=(H_x^\circ)^{L+1}C_x,&
 E_x^{[L]}&=R_x^{[L]}-\frac{\operatorname{Tr}R_x^{[L]}}{2q+1}I.
 \end{aligned}
 \tag{RMT19}
\]
The symbols \(B_x^\circ,H_x^\circ,\alpha_x\) label operators and
the scalar in this calculation; they leave every original
relation matrix \(B_N\), source metric \(M_x\), and coordinate
in (RMT1) in place.  Multiplication of the finite geometric sum
by \(I-H_x^\circ\) gives
\(C_x-C_x^{[L]}=R_x^{[L]}\) exactly.  Every factor is a real
function of the same \(M_0\)-self-adjoint \(D\), hence is
self-adjoint in \(M_x=M_0B_x^\circ\) as well.  Since the
trace of \(U_x-W_x\) vanishes, (RMT7) becomes
\[
 \mathcal B'(x)=
 \operatorname{Tr}\bigl(C_x^{[L]}(U_x-W_x)\bigr)
 +\operatorname{Tr}\bigl(E_x^{[L]}(U_x-W_x)\bigr).
\]
Apply Hilbert--Schmidt Cauchy--Schwarz in the original \(M_x\)
metric to the last term and integrate.  Define the specified
real numbers
\[
 \begin{split}
 j_L&=\int_0^1\operatorname{Tr}
                    \bigl(C_x^{[L]}(U_x-W_x)\bigr)\,dx,\\
 e_L&=\int_0^1
  \sqrt{\operatorname{Tr}((E_x^{[L]})^2)
                      \operatorname{Tr}((U_x-W_x)^2)}\,dx,\\
 \vartheta&=\frac{b_{2q+1}-b_1}{b_{2q+1}+b_1},\qquad
 K_D=\sum_{j=1}^{2q+1}
          \left(\frac{|b_j-1|}{\min\{1,b_j\}}\right)^2 .
 \end{split}
 \tag{RMT20}
\]
These are \(J_L,\epsilon_L,\vartheta,K_D\) of (ISM34)--(ISM37)
with \(j_L,e_L\) used to avoid collision with \(J_q\) in
(RMT13).  The eigenvalues of \(H_x^\circ\) have absolute
value at most \(\vartheta<1\); those of \(C_x\) are the
full \(c_j(x)\).  Thus
\(\operatorname{Tr}((E_x^{[L]})^2)
\leq\vartheta^{2L+2}K_D\), because subtracting the trace mean
subtracts a nonnegative square from the full sum of
squared residual eigenvalues.  The two positive
operators \(U_x,W_x\) dominate their range projections.
Those ranges intersect in dimension at least one, so
\(\operatorname{Tr}(U_xW_x)\geq1\), by positivity of the
trace product and the squared-length formula for the
product of the range projections.  Equation (RMT10) gives
\(\operatorname{Tr}((U_x-W_x)^2)\leq8q-6\).
The signed estimate and its computed residual are therefore
\[
 \begin{gathered}
 j_L-e_L\leq\mathcal B(1)-\mathcal B(0)\leq j_L+e_L,\\
 0\leq e_L\leq
       \vartheta^{L+1}\sqrt{K_D(8q-6)}
          \xrightarrow[L\to\infty]{}0 .
 \end{gathered}
 \tag{RMT21}
\]
Every integral is a continuous finite matrix expression.
This proves convergence for the fixed actual two source
forms; no uniformity in \(k\) is assumed or inferred.

Let \(B_0=\mathcal B(0)\) and retain \(H_0\) from (RMT18).
Intersecting the proved intervals, for every \(L\geq0\),
gives the propagated endpoint bounds
\[
 \begin{split}
 \mathcal B(1)&\geq
 \max\left\{0,\ B_0-J_q,\
       2a\log\cosh\left(\frac{\max\{0,H_0-I_q\}}2\right),\
       B_0+j_L-e_L\right\},\\
 \mathcal B(1)&\leq
 \min\left\{B_0+J_q,\
       2a\log\cosh\left(\frac{H_0+I_q}2\right),\
       B_0+j_L+e_L\right\}.
 \end{split}
 \tag{RMT22}
\]
Indeed each lower entry has separately been proved no
larger than the actual endpoint, and each upper entry no
smaller.  Taking their maximum and minimum therefore
preserves both inequalities and yields a nonempty
interval containing that endpoint.  A numerically
integrated center must retain its own integration
error before use; (RMT20)--(RMT22) state the exact
integrals.  The scalar source case has
\(H_x^\circ=E_x^{[L]}=0\) and \(j_L=e_L=0\),
so it also satisfies (RMT22) exactly.
